% !PRINTED 683
\weberpage{683}
%<<CONT
und wenn wir also annehmen, was offenbar keine Beschränkung
der Allgemeinheit ist, dass $\xi_1$ ein nicht verschwindender echter
Bruch sei, so widerspricht diese Formel der Annahme, dass

\[
\varepsilon_1, \varepsilon_2, \ldots, \varepsilon_{\nu-1}
\]

ein Fundamentalsystem sei, weil die Determinante $L_{\nu-1}$ verkleinert
 würde, wenn $\varepsilon_1$ durch $\varepsilon_0$ ersetzt wird.

Haben wir ein Fundamentalsystem, so können wir ein beliebiges
 System ganzer rationaler Zahlen $\xi_1, \xi_2, \ldots, \xi_{\nu-1}$ als
Exponentensystem annehmen und eine Einheit $\varepsilon$ mit diesen Exponenten
 bilden:

\[
\varepsilon = \varepsilon_1^{\xi_1} \varepsilon_2^{\xi_2} \ldots \varepsilon_{\nu-1}^{\xi_{\nu-1}}.
\]

Es bleibt also noch die Frage zu beantworten, inwieweit
eine Einheit durch das Exponentensystem bestimmt ist.

Nehmen wir an, es seien $\varepsilon', \varepsilon''$ zwei Einheiten mit demselben
Exponentensysteme, dann besteht das Exponentensystem der
Einheit $\varepsilon' : \varepsilon'' = \varrho$ aus lauter Nullen, und folglich sind die conjugirten
 Logarithmen von $\varrho$ alle gleich Null; $\varrho$ ist daher eine
ganze Zahl von der Eigenschaft, dass die absoluten Werthe aller
mit $\varrho$ conjugirten Zahlen gleich 1 sind, und daher, wie im
\S.\,175, 2. bewiesen ist, eine \gesperrt{Einheitswurzel}. Damit ist das
folgende Theorem bewiesen:

I. \gesperrt{Es giebt im Körper $\Omega$ ein System von $\nu - 1$ fundamentalen
 Einheiten}

\weq{5}{\varepsilon_1, \varepsilon_2, \ldots, \varepsilon_{\nu-1},}

\gesperrt{welches die Eigenschaft hat, dass in der Form}

\weq{6}{\varepsilon = \varrho\, \varepsilon_1^{\xi_1} \varepsilon_2^{\xi_2} \ldots \varepsilon_{\nu-1}^{\xi_{\nu-1}}}

\gesperrt{alle Einheiten des Körpers, jede nur einmal, enthalten sind,
wenn $\xi_1, \xi_2, \ldots, \xi_{\nu-1}$ alle ganzen rationalen Zahlen und $\varrho$ alle
in $\Omega$ vorhandenen Einheitswurzeln durchläuft.}

Es ist nun leicht, aus einem Fundamentalsysteme alle anderen
abzuleiten, indem man in (6) für die Exponenten $\nu - 1$ verschiedene
 Systeme ganzer Zahlen $\xi_{i,k}$ setzt, deren Determinante
$= \pm 1$ ist. Denn setzt man

\[
\varepsilon_i' = \varrho_i\, \varepsilon_1^{\xi_{1,i}} \varepsilon_2^{\xi_{2,i}} \ldots \varepsilon_{\nu-1}^{\xi_{\nu-1,i}},
\]

so folgt aus (3):

\[
L(\varepsilon_1', \varepsilon_2', \ldots, \varepsilon_{\nu-1}') = \pm L(\varepsilon_1, \varepsilon_2, \ldots, \varepsilon_{\nu-1}).
\]
